\subsection{RQ5. Is the interpretability of ML/DL modeled addressed in existing studies?}

%\song{give some motivation why interpretability is important}

%During the analysis of primary studies, we find 3 primary studies focusing on the interpretation of DL models for vulnerability detection ~\cite{3429444, li2021vulnerability, liu2021smart}. First, we explain possible strategies for model interpretation. Then, we discuss each paper in detail. 

%It is important to understand that there is a distinction between interpreting a DL model and selecting features.\song{can not understand} 

The goal of DL model interpretation is to identify the feature vectors that explain why a particular example is classified as a specific class. On the other hand, feature selection techniques such as Info Gain \cite{i32} and Chi-Squared \cite{i18} are primarily used to choose the most informative features for classification. The objective is to select a set of features that can better discriminate test samples from each other. It is worth noting that DL model interpretation is fundamentally different from generating adversarial examples \cite{i48, i50}. %During the analysis of primary studies, we find 3 primary studies focusing on the interpretation of DL models for vulnerability detection ~\cite{3429444, li2021vulnerability, liu2021smart}. 
Existing work on model interpretation can be classified into three broad categories including \textit{hidden neuron analysis}, \textit{model simulation}, and \textit{local interpretation}. 

\textbf{Hidden Neuron Analysis} is a strategy used to understand the features learned by a deep neural network by converting hidden neurons into a human-readable format. There are several methods used to achieve this. The first method \cite{i24, i42} involves measuring feature importance by analyzing gradients, which can lead to the vanishing gradients problem. The second method \cite{i43, i44, i47, i8} involves calculating feature importance using attention mechanisms, which identify important parts of the input. Model visualization is the third method, where the most informative hidden neurons are identified with respect to a specific class, such as using Global Max Pooling and ranked softmax weight in CNN models \cite{i52}. In the fourth method, intermediate results of CNNs are converted into the original input space to make it readable for humans, allowing for the observation of contours, feature colors, and local spaces of recognition across layers, such as with DeConvNet \cite{i49}. The fifth method connects the representation of concepts to hidden layers and measures the degree of matching \cite{i7}.

\noindent \textbf{Model Simulation} involves using more interpretable models such as linear regression, logistic regression, and decision tree to approximate and interpret models that are less explainable \cite{i3, i6, i11, i12, i14, i17, i23, i51}. For instance, it has been demonstrated that some neural networks can be explained and interpreted through local linear classifiers \cite{i11}. However, this comes with a trade-off of losing information and content that are learned by complex models when using linear models for interpretations.

\noindent \textbf{Local Interpretation.}
This approach is based on two main intuitions. The first intuition is that although the decision boundary of a global model may be complex, the decision boundary for a particular instance may be simple or linear in some cases. Two methods are used in this strategy, namely, \textit{sensitivity analysis} and \textit{local approximation}. The former aims to understand and characterize how a model's output is influenced by changing its inputs \cite{i16, i25, i29}. The latter, on the other hand, uses self-explanatory ML models to learn the importance of features corresponding to a specific sample \cite{i20, i22, i31}. For example, LIME \cite{i36} perturbs a given example to obtain a set of examples and feeds them to a linear regression model.

As mentioned earlier, in this survey, we identified three major studies focused on model interpretations~\cite{3429444, li2021vulnerability,liu2021smart}. In the first study \cite{3429444}, the authors utilized GNNexplainer, a model designed for interpreting vulnerable inputs at the statement level as part of the Program Dependence Graph (PDG) used as input representation. The PDG sub-graphs are believed to be associated with data and control dependencies among source code statements, which is why GNNexplainer \cite{ying2019gnnexplainer} was chosen to explain the rationale behind a model's decision. Specifically, when detecting vulnerabilities, the DL model's inputs and decisions are taken, along with the PDG of the sample, to identify PDG sub-graphs that explain the classification decision. The authors aimed to determine if the model's decision would be impacted when an edge is removed from the PDG sub-graph. If it does, the edge is considered informative.
The second study \cite{li2021vulnerability} presented a framework for a high-fidelity explanation of vulnerability detection by identifying small code tokens that have a significant impact on the model's output. In the first phase of the framework, the authors perturbed an example into variant examples near the decision boundary and searched for a sequence of important tokens. Specifically, they identified important tokens by measuring the impact of their perturbations on the prediction of the model for the variant examples. 
The third study \cite{liu2021smart} focuses on detecting vulnerabilities in Ethereum smart contracts. Conventional methods of detecting vulnerabilities in smart contracts rely on fixed expert rules, which have limited accuracy and scalability. Although recent deep learning approaches have improved the accuracy of detecting vulnerabilities, they fail to incorporate expert knowledge. The authors propose a method that combines deep learning with expert patterns in an explainable way. The method extracts expert patterns from the source code and casts them into a semantic graph to extract deep graph features. The global graph feature and local expert patterns are then fused to approach the final prediction, while their interpretable weights are also yielded. The experiments conducted on all available smart contracts on two platforms, Ethereum and VNT Chain, showed that their system significantly outperforms state-of-the-art methods. 

%\song{missing a summarization for RQ5} \nima{done}

\mybox{Answer to RQ5}{red!40}{gray!10}{
(1) There are three categories of model interpretation approaches, i.e., hidden neuron analysis, model simulation, and local interpretation. Hidden neuron analysis involves methods such as gradient analysis and attention mechanisms to understand the features learned by deep neural networks. The model simulation utilizes more interpretable models like linear regression and decision trees to approximate and interpret complex DL models. Local interpretation focuses on understanding the decision boundaries for specific instances and employs sensitivity analysis and local approximation methods.

(2) There are three specific studies in model interpretation for vulnerability detection~\cite{3429444, li2021vulnerability,liu2021smart}, highlighting the techniques and frameworks used in each study. These studies involve the interpretation of Program Dependence Graphs (PDGs), token impact analysis, and the fusion of deep learning with expert patterns. The findings and methodologies of these studies contribute to the field of explainable vulnerability detection in DL models.
}
\newpage